\phantomsection
\section*{ÖZET}
\addcontentsline{toc}{section}{ÖZET}
\vspace{.5cm}

Bu çalışmada, Arduino Uno tabanlı, iki eksenli hareket kabiliyetine sahip bir kızılötesi radar prototipi tasarlanmış ve gerçekleştirilmiştir. Proje, klasik RF radar sistemlerinin çalışma mantığını düşük maliyetli ve eğitim amaçlı bir gömülü sistem düzeneği üzerinde modellemeyi hedeflemektedir. Bu amaçla Sharp GP2Y0A02YK0F analog kızılötesi mesafe sensörü, yatay tarama için MG90S servo motor, dikey eksen hareketi için SG90S servo motor, yerel bilgilendirme için 16x2 I2C LCD ekran, hareket/titreşim gözlemi için MPU6050 ivmeölçer-jiroskop modülü ve veri kaydı için SPI haberleşmeli SD kart modülü kullanılmıştır.

Sistem, pan ekseninde çevreyi belirlenen açı aralığında taramakta, tilt ekseninde ise küçük açısal salınımlarla ölçüm doğrultusunu değiştirmektedir. Sensörden alınan analog mesafe verileri Arduino üzerinde işlenerek hedef var/yok kararı üretilmekte, anlık durum bilgileri LCD ekrana ve seri haberleşme üzerinden web arayüzüne aktarılmaktadır. Web arayüzü, Web Serial API ve HTML5 Canvas kullanılarak geliştirilmiş; pan açısı, tilt açısı, mesafe, hedef durumu, MPU6050 verileri, SD kart durumu ve Arduino besleme gerilimi gerçek zamanlı olarak görselleştirilmiştir. Böylece proje yalnızca sensör okuyan basit bir düzenek olmaktan çıkarılmış, gömülü yazılım, donanım entegrasyonu, veri kaydı ve kullanıcı arayüzü katmanlarını içeren bütünleşik bir radar simülasyon sistemi haline getirilmiştir.

Çalışma sürecinde servo motorların ani akım çekimi, ortak topraklama gereksinimi, I2C ve SPI hatlarının aynı sistemde kararlı çalıştırılması, sensör ölçümlerindeki gürültü ve seri haberleşmede eksik JSON paketlerinin ayıklanması gibi pratik mühendislik problemleri ele alınmıştır. Sonuç olarak sistem, eğitim ortamında radar prensiplerini, sensör füzyonunu ve gerçek zamanlı veri görselleştirmeyi gösterebilen uygulanabilir bir prototip olarak başarıyla çalıştırılmıştır.

\vspace{1cm}
\textbf{Anahtar Kelimeler:} Arduino Uno, Pan-Tilt Radar, Kızılötesi Mesafe Sensörü, MPU6050, SD Kart, Web Serial API

\newpage
\phantomsection
\section*{ABSTRACT}
\addcontentsline{toc}{section}{ABSTRACT}
\vspace{.5cm}

In this study, an Arduino Uno-based infrared radar prototype with two-axis motion capability was designed and implemented. The project aims to model the basic operating logic of radar systems on a low-cost educational embedded system platform. The hardware architecture includes a Sharp GP2Y0A02YK0F analog infrared distance sensor, an MG90S servo motor for horizontal scanning, an SG90S servo motor for vertical tilt motion, a 16x2 I2C LCD for local status output, an MPU6050 accelerometer-gyroscope module for motion and vibration monitoring, and an SPI-based SD card module for data logging.

The system scans the environment on the pan axis while applying small angular movements on the tilt axis. Analog distance measurements are processed on the Arduino to determine target presence, and telemetry data are transferred both to the LCD screen and to a browser-based interface through serial communication. The web interface was developed using the Web Serial API and HTML5 Canvas. It visualizes pan angle, tilt angle, distance, target status, MPU6050 readings, SD card status, and Arduino supply voltage in real time. Therefore, the project was developed as an integrated radar simulation system that combines embedded software, hardware integration, data logging, and user interface design.

During the implementation process, several practical engineering challenges were addressed, including sudden current draw of servo motors, common ground requirements, stable operation of I2C and SPI devices in the same system, noisy analog sensor readings, and filtering incomplete JSON packets in serial communication. As a result, the system was successfully operated as an applicable prototype that demonstrates radar principles, sensor-based measurement, and real-time data visualization in an educational environment.

\vspace{1cm}
\textbf{Keywords:} Arduino Uno, Pan-Tilt Radar, Infrared Distance Sensor, MPU6050, SD Card, Web Serial API
